% !TeX root = ../main.tex
% 第1.8节、第3.6节：中文摘要（篇幅以一页为宜，不出现图表、化学式和非公知术语）
\begin{abstractzh}
地球物理反演是地下介质结构成像与矿产资源勘探的核心技术手段。然而，传统基于最优化迭代的确定性反演方法通常面临计算开销大、对初始模型高度敏感以及强非线性条件下容易陷入局部极小值等瓶颈问题。近年来，数据驱动的深度学习方法在地球物理反演领域展现出显著优势，但在实际复杂地质场景中仍存在多尺度构造特征提取不足、泛化能力受限以及违背地球物理正演控制方程等突出挑战。

针对上述问题，本文围绕复杂地下介质的高精度智能反演展开系统研究，主要研究工作与创新性成果包括以下几个方面：

第一，提出了一种融合残差注意力机制的多尺度编解码网络结构。通过在编码端引入跨尺度空洞卷积金字塔与通道空间联合注意力模块，实现了对地下大尺度背景速度场与小尺度薄互层异常体的协同表征，有效缓解了深层网络下采样过程中的高频细节丢失问题。

第二，构建了嵌入正演物理算子的双驱动联合约束反演框架。将一维褶积与二维波动方程正演过程封装为可微物理层并嵌入训练损失函数，在不显著增加推理耗时的前提下，强制反演结果满足观测数据的动力学与运动学一致性，从而大幅提升了模型在少样本及跨工区迁移条件下的泛化能力。

第三，设计了基于蒙特卡洛随机失活与集成学习相结合的不确定性量化策略，在输出高分辨率地下物性参数估计的同时，定量给出各空间位置反演结果的置信区间与后验方差分布，为后续地质解释与钻井部署提供可靠的决策依据。

在多个合成基准数据集及某实际三维工区观测资料上的对比实验表明，与传统全波形反演方法及现有主流深度学习反演模型相比，本文方法在反演精度、抗噪鲁棒性以及计算效率方面均取得了显著提升，反演剖面的决定系数提高约 6.4\%，单道推理速度提升两个数量级以上，具有良好的工程应用价值。
\end{abstractzh}
